% Generated by scripts/materialize_unified_paper.py; do not edit.
\section{Standalone proof of the GNNW book lemma}

Date: 2026-08-12\\
Source being repaired: Gupta--Ndiaye--Norin--Wei, arXiv:2407.19026v1,
Lemmas 7--11 and Theorem 12.

\subsection{Claim boundary}

This note proves the combinatorial input used by the certified diagonal
Ramsey descent without invoking GNNW Lemma 11 or Theorem 12 as black
boxes. It retains their definitions of a candidate and of the rate
region, and it uses only elementary Ramsey recursion, convexity,
averaging, and induction.

Let a \textbf{candidate} be a pair of nonempty disjoint vertex sets
\((X,Y)\) in a red--blue complete graph. Its red cross-density is

\[
d(X,Y)=\frac{e_R(X,Y)}{|X||Y|}.
\]

It is \((k,\ell,t)\)-good if \(X\cup Y\) contains a red \(K_k\), or
\(X\) contains a blue \(K_t\), or \(Y\) contains a blue \(K_\ell\).

Let \(\mathcal R\) be the closure of the set of pairs
\((x,y)\in(0,1)^2\) for which, for all sufficiently large \(k+\ell\),

\[
R(k,\ell)\le x^{-k}y^{-\ell},
\]

and let \(\mathcal R_*\) be its interior.

\subsubsection{Theorem A (book induction)}

Suppose \(0<\mu_0,x_0,y_0,p<1\),

\[
x_0<p^{1/(1-\mu_0)}(1-\mu_0),
\qquad (x_0,y_0)\in\mathcal R_*.
\tag{U-B.A.1}\label{up:app-B:eq:a-1}
\]

There is \(L_0\) such that, for all positive integers \(k,\ell,t\) with
\(\ell\ge L_0\), every candidate of density at least \(p\) and
satisfying

\[
|X||Y|\ge x_0^{-k}y_0^{-\ell}\mu_0^{-t}
\tag{U-B.A.2}\label{up:app-B:eq:a-2}
\]

is \((k,\ell,t)\)-good.

\subsubsection{Corollary B (balanced-book bound)}

Suppose \(0<\mu,x,y,p<1\),

\[
x<p^{1/(1-\mu)}(1-\mu),
\qquad (x,y)\in\mathcal R_*.
\tag{U-B.B.1}\label{up:app-B:eq:b-1}
\]

For all sufficiently large \(\ell\), every red--blue coloring of \(K_N\)
of red density at least \(p\), where

\[
N\ge x^{-k/2}(\mu y)^{-\ell/2},
\tag{U-B.B.2}\label{up:app-B:eq:b-2}
\]

contains a red \(K_k\) or a blue \(K_\ell\).

The statements hold for all \(p,\mu\in(0,1)\); no relation \(p>\mu\) is
needed.

\subsection{Four elementary inputs}

\subsubsection{Lemma 1 (an interior point gives its own eventual bound)}

If \((u,v)\in\mathcal R_*\), then, for all sufficiently large \(a+b\),

\[
R(a,b)\le u^{-a}v^{-b}.
\tag{U-B.1.1}\label{up:app-B:eq:1-1}
\]

\textbf{Proof.} Choose \(\eta>0\) so that
\((u+2\eta,v+2\eta)\in\mathcal R\). Since \(\mathcal R\) is the closure
of the eventual-bound set, there is a point \((u',v')\) in that set with
\(u'>u\) and \(v'>v\). Its defining bound is stronger than \eqref{up:app-B:eq:1-1}.
\(\square\)

\subsubsection{Lemma 2 (density convexity)}

For every candidate,

\[
\sum_{v\in X}d(X,N_R(v)\cap Y)|N_R(v)\cap Y|
\ge e_R(X,Y)d(X,Y),
\tag{U-B.2.1}\label{up:app-B:eq:2-1}
\]

where an empty-neighborhood summand is interpreted as zero.

\textbf{Proof.} Put \(d_y=|N_R(y)\cap X|\). Multiplying the left side by
\(|X|\) gives \(\sum_{y\in Y}d_y^2\). Cauchy--Schwarz bounds this below
by \(|Y|^{-1}(\sum_y d_y)^2=e_R(X,Y)^2/|Y|\), which is \(|X|\) times the
right side. \(\square\)

\subsubsection{Lemma 3 (blue-book extraction)}

Let \(0<\nu<1\), and let \(b,k,m\) be positive integers with
\(m\ge5\nu^{-1}b^2\) and \(b\le m\). Suppose \(|X|\ge5m^2\) and at least
\(R(k,m)\) vertices \(v\in X\) have blue degree at least \(\nu|X|\)
inside \(X\). Then \(X\) contains a red \(K_k\), or it contains disjoint
sets \(S,T\) such that \(S\) is a blue clique, every edge from \(S\) to
\(T\) is blue, and

\[
|S|=b,\qquad |T|\ge\frac{\nu^b}{2}|X|.
\tag{U-B.3.1}\label{up:app-B:eq:3-1}
\]

\textbf{Proof.} The high-blue-degree vertices contain a red \(K_k\) or a
blue \(K_m\). In the latter case call its vertex set \(U\). Put
\(d_z=|N_B(z)\cap U|\) for \(z\in X\setminus U\) and let \(D\) be their
average. Counting ordered blue pairs from \(U\) to \(X\setminus U\)
gives

\[
D\ge\frac{m(\nu|X|-m)}{|X|-m}
=\nu m-\frac{(1-\nu)m^2}{|X|-m}
>\nu m-\frac14,
\tag{U-B.3.2}\label{up:app-B:eq:3-2}
\]

where \(|X|\ge5m^2\) and \(m\ge1\) were used in the last inequality. Let
\(q=\lfloor D\rfloor\). The discrete convexity of \(j\mapsto\binom jb\)
on nonnegative integers (with value zero for \(j<b\)) shows that a
uniform \(b\)-subset \(S\) of \(U\) satisfies

\[
\mathbb E|N_B(S)\cap(X\setminus U)|
=\binom mb^{-1}\sum_{z\in X\setminus U}\binom{d_z}b
\ge\binom mb^{-1}\binom qb|X\setminus U|.
\tag{U-B.3.3}\label{up:app-B:eq:3-3}
\]

Since \(q>D-1>\nu m-5/4\) and \(m\ge5\nu^{-1}b^2\), every \(0\le i<b\)
satisfies

\[
\frac{q-i}{m-i}
>\nu-\frac{(1-\nu)i+5/4}{m-i}
\ge\nu\left(1-\frac{b+1/4}{5b^2}\right).
\]

For the second inequality, the fraction being subtracted increases with
\(i\), and at \(i=b-1\) it is at most the claimed value because

\[
\frac{(1-\nu)(b-1)+5/4}{5b^2/\nu-(b-1)}
\le\frac{\nu(b+1/4)}{5b^2};
\]

after cross-multiplication, the difference between the right numerator
and the left numerator is \(\nu(b-1)(1-(b+1/4)/(5b^2))\ge0\).

(If \(q<b\), the last strict lower bound would be positive at \(i=q\),
an impossibility, so the binomial ratio is well-defined.) Hence

\[
\frac{\binom qb}{\binom mb}
\ge\nu^b\left(1-\frac{b+1/4}{5b^2}\right)^b
\ge\frac34\nu^b.
\tag{U-B.3.4}\label{up:app-B:eq:3-4}
\]

The last elementary bound is weakest at \(b=1\), where its left factor
is \(3/4\); for \(b\ge2\), Bernoulli's inequality gives at least
\(1-(b+1/4)/(5b)>3/4\). Finally \(|X\setminus U|\ge(4/5)|X|\), so
\eqref{up:app-B:eq:3-3}--\eqref{up:app-B:eq:3-4} exceed \(3\nu^b|X|/5>\nu^b|X|/2\). Some \(S\) therefore has
a common blue page \(T\) satisfying \eqref{up:app-B:eq:3-1}. \(\square\)

\subsubsection{Lemma 4 (the limit needed to choose the moment)}

For every \(p,\mu\in(0,1)\),

\[
\lim_{r\to\infty}
(p^{1/r}-\mu)^r(1-\mu)^{1-r}
=p^{1/(1-\mu)}(1-\mu),
\tag{U-B.4.1}\label{up:app-B:eq:4-1}
\]

where \(r\) runs through sufficiently large positive integers, for which
the base is positive.

\textbf{Proof.} Since \(p^{1/r}\to1>\mu\), positivity is eventual, and

\[
r\log\left(1+\frac{p^{1/r}-1}{1-\mu}\right)
\longrightarrow\frac{\log p}{1-\mu}.
\]

Exponentiation and the remaining factor \(1-\mu\) prove \eqref{up:app-B:eq:4-1}.
\(\square\)

\subsection{Proof of Theorem A}

The strict inequality in \eqref{up:app-B:eq:a-1} implies \(x_0+\mu_0<1\). By Lemma 4
choose an integer \(r>1\) such that

\[
x_0<(p^{1/r}-\mu_0)^r(1-\mu_0)^{1-r}.
\tag{U-B.5.1}\label{up:app-B:eq:5-1}
\]

After fixing \(r\), choose \(\varepsilon>0\) sufficiently small that

\[
\begin{gathered}
p\ge2\varepsilon,\qquad
(1+\varepsilon)(\mu_0+\varepsilon)\le1,\qquad
x_0+\mu_0+2\varepsilon<1,\\
(x_0+2\varepsilon,y_0+2\varepsilon)\in\mathcal R_*,\\
\mu_0+2\varepsilon
\le\frac{\mu_0+\varepsilon}
{(\mu_0+\varepsilon)+(x_0+\varepsilon)},\qquad
(p-\varepsilon)^{1/r}>\mu_0+3\varepsilon,\\
x_0\le
\big((p-\varepsilon)^{1/r}-\mu_0-3\varepsilon\big)^r
(1-\mu_0-\varepsilon)^{1-r}-\varepsilon.
\tag{U-B.5.2}\label{up:app-B:eq:5-2}
\end{gathered}
\]

All requirements follow by openness, continuity, \eqref{up:app-B:eq:5-1}, and the strict
inequality \(x_0+\mu_0<1\). Set

\[
x=x_0+\varepsilon,\qquad y=y_0+\varepsilon,\qquad
\mu=\mu_0+\varepsilon,\qquad \delta_n=\frac\varepsilon n.
\tag{U-B.5.3}\label{up:app-B:eq:5-3}
\]

We prove the following stronger assertion by induction on \(n=k+t\),
with \(\ell\ge L_0\) fixed. If

\[
d(X,Y)\ge p-\delta_n
\tag{U-B.5.4}\label{up:app-B:eq:5-4}
\]

and

\[
(d(X,Y)+\delta_n-p)^r|X||Y|
\ge x^{-k}y^{-\ell}\mu^{-t},
\tag{U-B.5.5}\label{up:app-B:eq:5-5}
\]

then \((X,Y)\) is \((k,\ell,t)\)-good. The cases \(k=1\) or \(t=1\) are
immediate.

First note that \eqref{up:app-B:eq:a-2} and \(d(X,Y)\ge p\) imply \eqref{up:app-B:eq:5-5}: because
\(x_0\le x/(1+\varepsilon)\), \(y_0\le y/(1+\varepsilon)\), and
\(\mu_0\le\mu/(1+\varepsilon)\), it is enough that

\[
\frac{\varepsilon^r}{n^r}(1+\varepsilon)^{n+\ell}\ge1.
\tag{U-B.5.6}\label{up:app-B:eq:5-6}
\]

A single sufficiently large \(L_0\) makes \eqref{up:app-B:eq:5-6} true for every
\(n\ge2\).

\subsubsection{\texorpdfstring{Regularization and the size of
\(X\)}{Regularization and the size of X}}

Delete from \(X\) any vertex having fewer than \((p-\delta_n)|Y|\) red
neighbors in \(Y\). Both the nonnegative excess

\[
e_R(X,Y)-(p-\delta_n)|X||Y|
\]

and its normalization by \(|X||Y|\) do not decrease, so the left side of
\eqref{up:app-B:eq:5-5} does not decrease. The positive right side prevents deletion of
all vertices. We may consequently assume

\[
d(X',Y)\ge p-\delta_n
\quad\text{for every nonempty }X'\subseteq X.
\tag{U-B.5.7}\label{up:app-B:eq:5-7}
\]

Lemma 1 and the interior condition in \eqref{up:app-B:eq:5-2} imply, for sufficiently
large \(L_0\), that if

\[
|Y|\ge(x+\varepsilon)^{-k}(y+\varepsilon)^{-\ell},
\]

then \(Y\) already contains a red \(K_k\) or a blue \(K_\ell\).
Otherwise, using \(d+\delta_n-p\le1\) in \eqref{up:app-B:eq:5-5},

\[
|X|
\ge\left(\frac{x+\varepsilon}{x}\right)^k
\left(\frac{y+\varepsilon}{y}\right)^\ell\mu^{-t}
\ge(1+\varepsilon)^{n+\ell}.
\tag{U-B.5.8}\label{up:app-B:eq:5-8}
\]

\subsubsection{The big-blue branch}

Put

\[
b=\left\lceil
\frac{2r\log n-r\log\varepsilon+\log2}{\log(1+\varepsilon)}
\right\rceil,\qquad
m=\lceil5\mu^{-1}b^2\rceil,\qquad w=R(k,m).
\tag{U-B.5.9}\label{up:app-B:eq:5-9}
\]

Here \(m=O_{r,\varepsilon}(\log^2n)\), \(m\ge b\), and
\(w=R(k,m)\le\binom{k+m-2}{m-1}\le(k+m)^m \le\exp(O_{r,\varepsilon}(\log^3n))\).
Since \eqref{up:app-B:eq:5-8} is exponential in \(n+\ell\), increasing \(L_0\) once makes,
uniformly for every \(n\),

\[
|X|\ge5m^2,\qquad
w\le\frac{\varepsilon(p-\varepsilon)}{n^3}|X|.
\tag{U-B.5.10}\label{up:app-B:eq:5-10}
\]

Let

\[
W=\{v\in X:|N_B(v)\cap X|\ge(\mu+\varepsilon)|X|\}.
\]

If \(|W|\ge w\), Lemma 3, with \(\nu=\mu+\varepsilon\), produces a red
\(K_k\) or a blue book with base \(S\) of size \(b\) and page \(T\) of
size at least \((\mu+\varepsilon)^b|X|/2\). If \(b\ge t\), its base
already contains a blue \(K_t\). Suppose \(b<t\). From \eqref{up:app-B:eq:5-7},
\(d(T,Y)\ge p-\delta_n\), and hence

\[
d(T,Y)+\delta_{n-b}-p
\ge\delta_{n-1}-\delta_n\ge\frac\varepsilon{n^2}.
\]

Using \eqref{up:app-B:eq:5-5}, \(d+\delta_n-p\le1\), and the definition of \(b\) gives

\[
\begin{aligned}
&(d(T,Y)+\delta_{n-b}-p)^r|T||Y|\\
&\quad\ge
\frac12\left(\frac\varepsilon{n^2}\right)^r
(\mu+\varepsilon)^b|X||Y|\\
&\quad\ge
\frac12\left(\frac\varepsilon{n^2}\right)^r
\left(\frac{\mu+\varepsilon}{\mu}\right)^b
x^{-k}y^{-\ell}\mu^{-t+b}\\
&\quad\ge x^{-k}y^{-\ell}\mu^{-t+b}.
\end{aligned}
\tag{U-B.5.11}\label{up:app-B:eq:5-11}
\]

Induction applied to \((T,Y)\) with parameter \(t-b\) finishes this
branch: a blue \(K_{t-b}\) in \(T\) joins the blue base \(S\).

We may therefore assume

\[
|W|<w.
\tag{U-B.5.12}\label{up:app-B:eq:5-12}
\]

\subsubsection{The red-step/density-increment branch}

Lemma 2 and \eqref{up:app-B:eq:5-10}--\eqref{up:app-B:eq:5-12} show that the contribution from \(W\) is at
most

\[
w|Y|\le\frac\varepsilon{n^3}e_R(X,Y).
\]

Consequently some \(v\in X\setminus W\) satisfies

\[
d(X,N_R(v)\cap Y)
\ge d(X,Y)-\frac\varepsilon{n^3}.
\tag{U-B.5.13}\label{up:app-B:eq:5-13}
\]

Write

\[
X_R=N_R(v)\cap X,\quad X_B=N_B(v)\cap X,\quad
Y'=N_R(v)\cap Y,\quad p'=p-\delta_{n-1},
\]

and

\[
\alpha=d(X,Y')-p',\quad
\alpha_R=d(X_R,Y')-p',\quad
\alpha_B=d(X_B,Y')-p'.
\]

An empty \(X_R\) or \(X_B\) is assigned zero weighted contribution and
its induction branch is skipped; its corresponding \(\alpha_R\) or
\(\alpha_B\) need not be defined. By \eqref{up:app-B:eq:5-7},
\(|Y'|\ge(p-\varepsilon)|Y|\). Also, \eqref{up:app-B:eq:5-4} and \eqref{up:app-B:eq:5-13} give

\[
\alpha\ge
\delta_{n-1}-\delta_n-\frac\varepsilon{n^3}
\ge\frac\varepsilon{n^2}>0.
\tag{U-B.5.14}\label{up:app-B:eq:5-14}
\]

Since all edges from \(v\) to \(Y'\) are red,

\[
\frac{\alpha_R}{\alpha}\frac{|X_R|}{|X|}
+\frac{\alpha_B}{\alpha}\frac{|X_B|}{|X|}
+\frac1{\alpha|X|}\ge1.
\tag{U-B.5.15}\label{up:app-B:eq:5-15}
\]

If \(\alpha_R\ge0\) and

\[
\alpha_R^r|X_R|
\ge\frac{x\alpha^r|X|}{p-\varepsilon},
\tag{U-B.5.16}\label{up:app-B:eq:5-16}
\]

then, using \(|Y'|\ge(p-\varepsilon)|Y|\), \eqref{up:app-B:eq:5-13}, and

\[
d(X,Y')-p'\ge d(X,Y)+\delta_n-p,
\]

condition \eqref{up:app-B:eq:5-16} and \eqref{up:app-B:eq:5-5} show that \((X_R,Y')\) satisfies the
inductive moment inequality for \((k-1,\ell,t)\). Its goodness implies
that of \((X,Y)\) because \(v\) is red to \(X_R\cup Y'\). The analogous
condition, including its necessary sign hypothesis, is

\[
\alpha_B\ge0,
\qquad
\alpha_B^r|X_B|
\ge\frac{\mu\alpha^r|X|}{p-\varepsilon}
\tag{U-B.5.17}\label{up:app-B:eq:5-17}
\]

gives the inductive branch \((k,\ell,t-1)\), since \(v\) is blue to
\(X_B\).

Assume both branches fail. Set

\[
a=|X_R|/|X|,\qquad c=|X_B|/|X|,\qquad q=1-1/r.
\]

Using the strict failures of \eqref{up:app-B:eq:5-16}--\eqref{up:app-B:eq:5-17} in \eqref{up:app-B:eq:5-15} gives

\[
x^{1/r}a^q+\mu^{1/r}c^q
+\frac{(p-\varepsilon)^{1/r}}{\alpha|X|}
>(p-\varepsilon)^{1/r}.
\tag{U-B.5.18}\label{up:app-B:eq:5-18}
\]

Now \(a+c=1-|X|^{-1}<1\) and \(c\le\mu+\varepsilon\). The function

\[
f(z)=x^{1/r}(1-z)^q+\mu^{1/r}z^q
\]

is increasing on \(0\le z\le\mu/(\mu+x)\). The special smallness
condition in \eqref{up:app-B:eq:5-2} is exactly

\[
\mu+\varepsilon\le\frac\mu{\mu+x}.
\]

Therefore

\[
x^{1/r}a^q+\mu^{1/r}c^q
\le f(c)\le f(\mu+\varepsilon)
\le x^{1/r}(1-\mu)^q+\mu+\varepsilon.
\tag{U-B.5.19}\label{up:app-B:eq:5-19}
\]

Equations \eqref{up:app-B:eq:5-8} and \eqref{up:app-B:eq:5-14}, after one final uniform increase of \(L_0\),
give

\[
\frac{(p-\varepsilon)^{1/r}}{\alpha|X|}
\le\frac{n^2}{\varepsilon(1+\varepsilon)^{n+\ell}}
\le\varepsilon.
\tag{U-B.5.20}\label{up:app-B:eq:5-20}
\]

Substituting \eqref{up:app-B:eq:5-19}--\eqref{up:app-B:eq:5-20} into the strict inequality \eqref{up:app-B:eq:5-18} yields

\[
x^{1/r}(1-\mu)^{1-1/r}+\mu+2\varepsilon
>(p-\varepsilon)^{1/r}.
\]

The positivity condition in \eqref{up:app-B:eq:5-2} allows the rearrangement

\[
x>
\big((p-\varepsilon)^{1/r}-\mu-2\varepsilon\big)^r
(1-\mu)^{1-r}.
\tag{U-B.5.21}\label{up:app-B:eq:5-21}
\]

But \(x=x_0+\varepsilon\) and \(\mu=\mu_0+\varepsilon\), so \eqref{up:app-B:eq:5-21}
strictly contradicts the final inequality of \eqref{up:app-B:eq:5-2}. The induction, and
hence Theorem A, is complete. \(\square\)

\subsection{Proof of Corollary B}

By openness choose \(y_0>y\) with \((x,y_0)\in\mathcal R_*\). Take a
uniformly random equitable bipartition of the \(N\) vertices. Each edge
has the same probability of crossing, so some partition \((X,Y)\) has
cross-red density at least the global red density \(p\).

For all sufficiently large \(\ell\),

\[
(y_0/y)^{\ell/2}\ge3.
\]

Thus \eqref{up:app-B:eq:b-2} makes both equitable parts at least
\(x^{-k/2}(\mu y_0)^{-\ell/2}\), and hence

\[
|X||Y|\ge x^{-k}y_0^{-\ell}\mu^{-\ell}.
\]

Apply Theorem A with \(t=\ell\). A red \(K_k\) in \(X\cup Y\), a blue
\(K_\ell\) in \(X\), or a blue \(K_\ell\) in \(Y\) gives the desired
conclusion. \(\square\)

\subsection{Source defects neutralized by this proof}

The standalone proof explicitly repairs all of the following
source-level issues rather than silently importing them:

\begin{enumerate}
\def\labelenumi{\arabic{enumi}.}
\tightlist
\item
  Lemma 10 does not require \(p>\mu\).
\item
  The moment order \(r\) is chosen before \(\varepsilon\).
\item
  The northeast rate point is chosen in \(\mathcal R_*\), and Lemma 1
  supplies the exact eventual bound used later.
\item
  The printed comparison for \(y_0\) uses \(x\) in place of \(y\).
\item
  The case \(b\ge t\) and empty \(X_R,X_B\) are treated separately.
\item
  The normalized sizes satisfy \(a+c<1\), not the dimensionally false
  printed inequality.
\item
  The concavity calculation supplies an upper bound in \eqref{up:app-B:eq:5-19}, and the
  strictness in \eqref{up:app-B:eq:5-18} is preserved to obtain a genuine contradiction.
\end{enumerate}

The only non-symbolic dependency left in the diagonal descent is the
stated semantics of the Arb interval implementation used to certify its
numerical inequalities. A post-repair adversarial replay of this standalone
proof and a generated-copy fidelity audit have both passed. Independent
external human review and public archival remain pending.
